All sensitivity runs use 40\% noise and the same 5 seeds; the main-setting rows are the main-grid runs (not re-run). The sweeps are evaluated on the test split and are \emph{not} used to revise the main results.

\begin{table}[t]\centering\caption{Sensitivity at 40\% noise (mean $\pm$ std over 5 seeds). CE reference: test\_acc 0.709, mem\_rate 1.000.}\label{tab:sens}
\resizebox{\columnwidth}{!}{\begin{tabular}{llccc}
\toprule
System & Setting & test\_acc & mem\_rate & mem\_gap \\
\midrule
Mixup & $\alpha$=0.2 & 0.709$\pm$0.037 & 0.984$\pm$0.010 & 0.976$\pm$0.015 \\
Mixup & $\alpha$=0.5 & 0.735$\pm$0.022 & 0.958$\pm$0.020 & 0.928$\pm$0.033 \\
Mixup & $\alpha$=1 (main) & 0.762$\pm$0.030 & 0.892$\pm$0.006 & 0.819$\pm$0.020 \\
Mixup & $\alpha$=2 & 0.794$\pm$0.026 & 0.718$\pm$0.013 & 0.486$\pm$0.029 \\
Mixup & $\alpha$=4 & 0.849$\pm$0.038 & 0.426$\pm$0.071 & -0.085$\pm$0.156 \\
\midrule
LS & $\epsilon$=0.05 & 0.662$\pm$0.042 & 1.000$\pm$0.000 & 1.000$\pm$0.000 \\
LS & $\epsilon$=0.1 (main) & 0.668$\pm$0.028 & 0.999$\pm$0.002 & 0.999$\pm$0.002 \\
LS & $\epsilon$=0.2 & 0.676$\pm$0.034 & 0.999$\pm$0.001 & 0.999$\pm$0.002 \\
LS & $\epsilon$=0.4 & 0.682$\pm$0.036 & 1.000$\pm$0.000 & 1.000$\pm$0.000 \\
LS & $\epsilon$=0.6 & 0.693$\pm$0.030 & 1.000$\pm$0.001 & 0.999$\pm$0.002 \\
\midrule
Small-loss & rate=0.1 & 0.733$\pm$0.029 & 0.875$\pm$0.005 & 0.778$\pm$0.009 \\
Small-loss & rate=0.2 & 0.774$\pm$0.023 & 0.685$\pm$0.015 & 0.418$\pm$0.027 \\
Small-loss & rate=0.3 & 0.828$\pm$0.016 & 0.509$\pm$0.015 & 0.082$\pm$0.034 \\
Small-loss & rate=0.4 (main) & 0.896$\pm$0.018 & 0.280$\pm$0.037 & -0.370$\pm$0.084 \\
Small-loss & rate=0.5 & 0.918$\pm$0.016 & 0.162$\pm$0.010 & -0.612$\pm$0.026 \\
Small-loss & rate=0.6 & 0.921$\pm$0.019 & 0.116$\pm$0.021 & -0.702$\pm$0.056 \\
\midrule
Small-loss & no warm-up/ramp, rate=0.4 & 0.924$\pm$0.019 & 0.146$\pm$0.040 & -0.649$\pm$0.091 \\
\bottomrule
\end{tabular}
}
\end{table}

\begin{table}[t]\centering\caption{No warm-up/ramp minus main-grid small-loss at 40\% noise, paired over 5 seeds.}\label{tab:warm}
\resizebox{\columnwidth}{!}{\begin{tabular}{lrr}
\toprule
Metric & paired $\Delta$ (no warm-up $-$ main) & $p$ \\
\midrule
test\_acc & +0.028 & 0.0030 \\
mem\_rate & -0.133 & $<$0.0001 \\
\bottomrule
\end{tabular}
}
\end{table}

\begin{figure}[t]\centering\includegraphics[width=\columnwidth]{sweeps.pdf}
\caption{Test accuracy at 40\% noise against the method hyperparameter (mean $\pm$ std, 5 seeds). Ring: main setting; dashed line: CE.}\label{fig:sweeps}\end{figure}

\textbf{H5 (``under-estimating the noise rate loses most of the gain''): partly supported, and the sign of the error matters.} Under-estimating the noise rate costs accuracy monotonically; most of the gain over CE is lost at assumed rates 0.1 and 0.2, but not at 0.3: assumed rates 0.1, 0.2, 0.3 give 0.733, 0.774, 0.828 against 0.896 at the true rate 0.4, and mem\_rate rises from 0.280 to 0.875 at 0.1, close to CE (0.709/1.000). Over-estimating does \emph{not} hurt here: rates 0.5 and 0.6 give 0.918 and 0.921, higher than the true rate, within roughly one standard deviation of each other. Discarding more than the true fraction removes some clean examples but, on a dataset this redundant (about 1,000 examples of 10 well-separated classes), still leaves enough clean ones. This is an observation for this dataset only; it would not be expected to hold when classes are rare or hard. We did not sweep rates above 0.6.

\textbf{Warm-up and ramp.} Removing the 10-epoch warm-up and the 10-epoch ramp (selection at the full rate from epoch 0) gives 0.924 against 0.896 for the main setting (Table~\ref{tab:sens}) and mem\_rate 0.146 against 0.280. The main schedule was fixed before the runs. The comparison is paired over the same five seeds and noise draws (Table~\ref{tab:warm}; uncorrected $t$-test, evaluated on the test split), and the direction is consistent with memorisation having started during the warm-up. This also means the main-grid small-loss numbers are conservative with respect to this component, not tuned.

\textbf{Mixup strength.} Test accuracy rises monotonically with $\alpha$ (0.709, 0.735, 0.762, 0.794, 0.849 for $\alpha=0.2,0.5,1,2,4$) while mem\_rate falls from 0.984 to 0.426; at $\alpha=4$ the mem\_gap is $-0.085\pm0.156$, i.e.\ roughly balanced. The mixup number in the main grid ($\alpha{=}1$) is therefore a weak setting for mixup on this task; the main comparison was not re-done with $\alpha{=}4$, since choosing it after seeing test results would be tuning on the test split. At $\alpha{=}4$ mixup (0.849) is still below small-loss at the true rate (0.896) but this is a different comparison from the pre-registered one, and the sweep did not include the matched sweep over the small-loss schedule ($T_w$, $T_r$) that would be needed for a fair tuned comparison.

\textbf{Label smoothing strength.} Accuracy changes only from 0.662 to 0.693 over a 12-fold range of $\epsilon$ with mem\_rate pinned at 1.0, so the smoothing target does not prevent memorisation in this regime; this is in line with the report that its benefit vanishes at high noise~\cite{wei2021smooth}, though our setting differs.
